\section{Method}
\label{sec:method}

\subsection{Age-aligned windows}

For an engine at cut $c$, let $t_0=c-1$ be the index of the most recent cycle.
For a window lag $\Delta$ and offset $j$, the window ends at $t_j=t_0-j\Delta$.
The backbone prediction for that window is $p_j$. Because the window is older
than the cut, its prediction refers to an earlier cycle. We align it to the cut:
\begin{equation}
a_j = p_j + j\Delta .
\label{eq:align}
\end{equation}
The aligned predictions $\{a_j\}_{j=0}^{K-1}$ are the inputs to the aggregator.
For $K=1$, $a_0=p_0=K1$; for $K=3$ and $\Delta=5$ we use offsets $j=0,1,2$.

\subsection{Features}

The aggregator uses eleven features: the single-window prediction $p_0$, the
aligned median $p_{\mathrm{med}}$, the aligned mean $p_{\mathrm{mean}}$, the
aligned standard deviation $p_{\mathrm{std}}$, the median absolute deviation
(MAD), the shift $p_{\mathrm{med}}-p_0$, the regime distance, the cap indicator,
the current cut normalized by 250, and the observed history length normalized by
70. The lifetime fraction $c/n$ is deliberately excluded: the true lifetime $n$
is not observable at test time. The cap indicator is computed from the model
outputs ($p_0\ge 124.5$ or $p_{\mathrm{med}}\ge 124.5$), not from the true RUL.

\subsection{Learned gate}

LTAG is a two-layer MLP with GELU activation and a sigmoid output:
\begin{equation}
g(x)=\sigma\!\left(W_2\,\mathrm{GELU}(W_1 x+b_1)+b_2\right),
\end{equation}
with the last layer initialized to zero, so the gate starts at $g=0.5$. The
prediction is
\begin{equation}
\hat y = g\,p_{\mathrm{med}} + (1-g)\,p_0 .
\label{eq:gate}
\end{equation}
The gate is trained with a censored hinge loss on the capped target,
\begin{equation}
\ell(\hat y,y)=
\begin{cases}
\max(0,\hat y-y), & y<125,\\
\max(0,\hat y-125), & y=125,
\end{cases}
\end{equation}
which avoids penalizing a prediction above the cap when the true RUL is capped.
Training uses AdamW, learning rate $10^{-3}$, weight decay $10^{-3}$, 300 epochs,
and five engine-disjoint folds (GroupKFold by subset and engine). Validation
engines are not used in any fold.

\subsection{Calibration baselines}

We compare two post-hoc calibrators on the same features. K1-Calib is a residual
MLP that predicts $\hat y=p_0+r(x)$, where $r$ is a two-layer MLP with zero-init
last layer and the input contains only $p_0$ and context (regime distance, cap,
cut, history). LS-ctx is a ridge regression on
$[p_0,p_{\mathrm{med}},\text{context}]$. Both are trained on the training engines
and applied to validation engines; neither sees test labels.

The distinction between LTAG and the calibrators is deliberate. LTAG learns a
\emph{gate} between two existing predictions; K1-Calib learns an additive
correction to $p_0$; LS-ctx learns a linear map. If a calibrator matches LTAG,
the gate is not necessary. Section~\ref{sec:cutwise} shows that the answer
depends on the cut.
